\documentclass[12pt]{scrartcl}

\usepackage[utf8]{inputenc}
\usepackage[T1]{fontenc}
\usepackage[english]{babel}
\usepackage{lmodern}
\usepackage{graphicx}
\usepackage[pdftex,hyperref,dvipsnames]{xcolor}
\usepackage{listings}
\usepackage[a4paper,lmargin={2cm},rmargin={2cm},tmargin={3.5cm},bmargin = {2.5cm},headheight = {4cm}]{geometry}
\usepackage{amsmath,amssymb,amstext,amsthm}
\usepackage[lined,algonl,boxed]{algorithm2e}
\usepackage{tikz}
\usepackage{hyperref}
\usepackage{url}
\usepackage[inline]{enumitem}
\usepackage[headsepline]{scrlayer-scrpage}
\usepackage{pgfplots}
\usepackage{caption}
\usepackage{array}
\usepackage{ifthen}
\usepackage{multirow}
\usepackage{yfonts}

\pagestyle{scrheadings}
\usetikzlibrary{automata,positioning}


% Überschreibt enumerate Befehl, sodass 1. Ebene Items mit
\renewcommand{\theenumi}{(\alph{enumi})}
% (a), (b), etc. nummeriert werden.
\renewcommand{\labelenumi}{\text{\theenumi}}

\newcolumntype{L}[1]{>{\raggedright\let\newline\\\arraybackslash\hspace{0pt}}m{#1}}
\newcolumntype{C}[1]{>{\centering\let\newline\\\arraybackslash\hspace{0pt}}m{#1}}
\newcolumntype{R}[1]{>{\raggedleft\let\newline\\\arraybackslash\hspace{0pt}}m{#1}}


% \setcounter{countername}{number}: Legt den Wert des Counters fest
% \stepcounter{countername}: Erhöht den Wert des Counters um 1.
\newcounter{sheetnr}
\newcounter{exnum}
\newcounter{subexnum}
\newcounter{subsubexnum}

\newcommand{\pic}[2][0.7]{
    \begin{figure}[h!]
        \centering
        \includegraphics[scale=#1]{#2}
    \end{figure}
}

\newcommand{\picwithcap}[4][0.7]{
    \begin{figure}[h!]
        \centering
        \includegraphics[scale=#1]{#2}
        \caption{#3}
        \label{#4}
    \end{figure}
}

\ohead{
    Felix R\"ohr 3567731 st176436@stud.uni-stuttgart.de}
% \chead{} kann mittleren Kopfzeilen Teil sezten
% \ihead: Setzt linken Teil der Kopfzeile mit
% Modulnamen, Semester und Übungsblattnummer
\ihead{Proposal Bachelor thesis}

\title{{\small Bachelor thesis}\\ Ballot Validity Proofs, An analysis of instantiations of ZK-SNARKs for Ensuring Ballot Validity}

\begin{document}

\maketitle

\begin{itemize}
    \item Student: Felix Röhr
    \item Examiner: Prof. Dr. Ralf Küsters
    \item Supervisor : Nicolas Huber
    \item Start: October 2024, End: April 2025
\end{itemize}

\section{Thesis topic}
In today's world, the necessity of secure e-voting schemes is undisputed. 
For this purpose, Zero Knowledge Proofs (ZKPs) are used to allow users to prove the validity of their votes without revealing them. 
One class of ZKPs that can be applied to various e-voting systems are General Purpose Zero Knowledge proofs (GPZKPs). 
In practice, Zero-knowledge Succinct Non-interactive Arguments of Knowledge (zk-SNARKs) are often used to create GPZKPs.

The feasability of GPZKPs for showing ballot validity in voting systems with arbitrary election types and ballot formats using Groth16-SNARKs \cite{Groth16} has been shown by Huber et. al \cite{BallotValidity}.

As examples, the authors examined the different election types single vote, multi-vote, "Line-Vote", "Multi-Vote with Rules", the ranked election types Pointlist-Borda and Borda Tournament Style, and Cordocet methods.

Ballot validity for an election type encompasses two statements: First, that a ballot $b$ is chosen from the set of allowed ballots $C$ for that election type, and second, that the public encrypted ballot $c$ is actually the encryption $Enc(b)$ of the ballot $b$. 
The indicator function $f_R(x,w)$ stating the truth value of the conjunction of these two statements $\mathfrak{C}_{\text{Voting}}$ and $\mathfrak{C}_{\text{Enc}}$ is represented as an arithmetic circuit, where inputs are elements of a scalar field $\mathbb{F}$.
Finally, this circuit gets further transformed into a into a Rank-1 constraint system (R1CS) and from there into a quadratic arithmetic program (QAP) which is used by Groth16 to specify $f_R$ \cite{BallotValidity,lambdaclass}.
Groth16 SNARKs consist of three group elements over a base field. The size of this base field is determined by the elliptic curve the SNARK is instantiated with \cite{BallotValidity}.

Huber et. al used libsnark as an implementation of Groth16 SNARKs, which uses the elliptic curve BN254 that provides a base field of size $\approx 2^{254}$. Furthermore, Exponential ElGamal Encryption was used for $Enc$ \cite{BallotValidity}.

This thesis will focus on implementing the ballot validity proofs for the different election types covered in \cite{BallotValidity} in Circom \cite{circom}. Circom together with the SNARK-backend snarkjs provides another implementation of Groth16 SNARKs. Building on this, the thesis will compare the performance of the construction of ballot validity proofs in Circom using the elliptic curve BN254 for the different election types.
Furthermore, the performance will be compared to that of libsnark under the same circumstances.

Circom is advantageous when compared to libsnark, since it provides a high-level language to define circuits making circuit construction simpler and faster.
Moreover, Circom supports two different elliptic curves, namely BN254 which is also supported by libsnark and BLS12-381 \cite{circom}.

Therefore, if time permits, this thesis can also include an implementation of the ballot validity proofs from \cite{BallotValidity} using the BLS12-381 elliptic curve. 
Then, the comparative analysis of libsnark and circom using the BN254 curve could be extended to include a comparison with the BLS12-381 curve.

Additionally, the work for this thesis can include an implementation of the ballot validity proofs for the aformentioned election types in another Groth16-SNARK backend that supports BN254 such as gnark \cite{gnark}.
The results of this can then be included in the comparative analysis.

\section{Tasks}
\subsection{Mandatory}
This thesis includes:
\begin{itemize}
    \item Research into:
    \begin{itemize}
        \item Representing statements as arithmetic circuits
        \item Translating arithmetic circuits into a format understandable by Circom \cite{lambdaclass, circom, MoonMath}
        \item Groth16-SNARKs \cite{MoonMath}
        \item Voting systems, election types and ballot validity
        \item Circom \cite{circom}
        \item BN254 elliptic curve
    \end{itemize} 
    \item Transcribing the arithmetic circuits representing ballot validity for the different election types \cite{BallotSnarks} mentioned above into formats suitable for circom
    \item Performance comparison for the construction of proofs for ballot validity for the different election types in circom and between circom and libsnark
\end{itemize}

\subsection{Optional}
This thesis can include:
\begin{itemize}
    \item Research into:
    \begin{itemize}
        \item BLS12-381 elliptic curve
        \item gnark \cite{gnark}
    \end{itemize} 
    \item Transcribing the arithmetic circuits representing ballot validity for the different election types \cite{BallotSnarks} mentioned above into formats suitable for gnark
    \item Performance comparison for the construction of proofs for ballot validity for the different election types in gnark and between gnark and the other SNARK-backends libsnark and circom (using the BN254 curve)
    \item Performance comparison for the construction of proofs for ballot validity for the different election types in circom and ganrk unsing the BLS12-381 curve
\end{itemize}

\bibliographystyle{plain}
\bibliography{refs}

% - groth16 SNARKs 
% - optional: andere SNARKs
% 
% - obligatorisch: Circom mit einer vorgegebenen Kurve (BN254).
% 
% - Verschiedenen eliptische Kurven, basieren auf versch. Gruppenpaarungen (diese vergl.)
% - Base field vs. scalar field
% - andere Impl., die andere elliptische Kurven erlauben (nicht libsnark)
% 
% - ballot validity:
%     - in elliptischen Kurven rechnen
% 
% - e.g., libsnark vs. circom vs. gnark (haben evtl. eigene Backends bzw. ) (optional)
% 
% - optional: versch. Kurven in circom vgl.
% 
% - optional: auch noch Plonk mit Groth16 vgl.
% 
% Nächste Schritte:
% - Einarbeiten in Circom (Achtung: Keine Operationen außerhalb des vorgegebenen Scalar fields)
% - Proposal noch mal überarbeiten
% - MoonMathManual: Montgomery Kurven sind relevant, der Rest von elliptischen Kurven nicht so wirklich
% (Pairing Alg. eher nicht so relevant)
% 
% - MoonMathManual: Kapitel 6-8 v. a. relevant

\end{document}